\section{الفصل~\ref{sec:dofa}. التعلّم مع \nt{DOPA}}

آليات المشبك (الدفعات، وكاشف التزامن، والكالسيوم، ونوافذ اللدونة) وسلوك الدوبامين خطأً للتنبؤ مقيسة مباشرة؛ وأن القواعد تقود إلى التدرج، وما ثمن البثّ، مبرهنتان؛ ونتيجة تعلّم الاختيار حساب بنموذج الكتاب؛ والدارة التي تحسب الخطأ فرضية.

{\footnotesize
\setlength{\tabcolsep}{3pt}\renewcommand{\arraystretch}{1.15}
\begin{longtable}{>{\raggedright}p{0.5cm}>{\raggedright}p{4.0cm}>{\raggedright}p{5.6cm}>{\raggedright}p{2.3cm}>{\raggedright\arraybackslash}p{1.7cm}}
\toprule
م & القضية & التجربة وما قيس & المصدر & الحالة \\
\midrule\endfirsthead
\toprule
م & القضية & التجربة وما قيس & المصدر & الحالة \\
\midrule\endhead
1 & لم يُعثر في الدماغ على الانتشار العكسي للخطأ بالشكل الذي هو عليه في الشبكات الاصطناعية & لا توجد تجربة مباشرة: هذه قضية عن غياب. تحلل المراجعات ما تتطلبه الخوارزمية (وصلات راجعة تكرر الأمامية، وطور مستقل) وما اقتُرح لها من تقريبات بيولوجية. & \cite{Lillicrap2020, WhittingtonBogacz2019} & فرضية \\
2 & يتفكك الوزن إلى عدد مواقع التحرير واحتمال التحرير والاستجابة للدفعة: $w=N_s\,p\,A_1$~\eqref{eq:quantal} & المشبك العصبي العضلي عند الضفدع مع تحرير مخفَّض: تتقلب استجابة المتلقي بدرجات هي مضاعفات الاستجابة المصغّرة التلقائية لدفعة واحدة؛ وعدد الدفعات لكل نبضة يخضع لقانون بواسون. & \cite{delCastillo1954} & مُقاس \\
3 & مستقبِل المتلقي كاشف تزامن؛ والكالسيوم يطلق تغيّر الوزن & تثبيت الجهد الرقعي (\foreignlanguage{english}{patch clamp}) لعصبونات الفأر: أيونات \foreignlanguage{english}{Mg$^{2+}$} تسدّ القناة التي يفتحها الغلوتامات عند كمون الراحة وتخرج منها عند إزالة الاستقطاب. شرائح الحُصين: خالب الكالسيوم في المتلقي يحجب التقوية الطويلة الأمد، وتحرير الكالسيوم بالضوء داخل الخلية يقوّي النقل بنفسه. & \cite{Nowak1984, Malenka1988calcium} & مُقاس \\
4 & القرار ينفذه أي جانب: يكبر $A_1$ عند المتلقي، ويتغير $p$ عند المرسِل & الغلوتامات المحرَّر بالضوء فوق شويكة واحدة يسبب نموًا سريعًا وثابتًا للشويكة وللتيار عبر مستقبلاتها؛ وترافق التقويةَ إدماجُ مستقبلات \foreignlanguage{english}{AMPA}. إزالة استقطاب المتلقي تكبت مؤقتًا التحرير عند المرسِل؛ وينقل الأثرَ القنبينوداتُ الداخلية التي تعبر الشق عائدة (أُظهر ذلك على مداخل \nt{GABA}). والاستنزاف القصير الأمد يعتمد على احتمال التحرير عند المرسِل. & \cite{Matsuzaki2004, Malinow2002, WilsonNicoll2001, Tsodyks1997} & مُقاس \\
5 & يقوى المشبك حين يكون المرسِل والمتلقي نشطين معًا~\eqref{eq:hebb} & أرنب مخدَّر: بعد سلاسل قصيرة من النبضات على مسار الدخل إلى الحُصين تبقى الاستجابة مقوّاة مدة من $30$ دقيقة إلى $10$ ساعات. وفي شريحة الحُصين لا تقوّي نبضات المتلقي إلا المشابك النشطة في اللحظة نفسها. & \cite{Bliss1973LTP, Kelso1986hebb} & مُقاس \\
6 & القاعدة~\eqref{eq:oja} تجد المتجه الذاتي الرئيسي لارتباطات المداخل (القضية~\ref{prop:oja}) & مبرهنة؛ والبرهان في الفصل. لا توجد تجربة تعلّم فيها عصبون المكوّن الرئيسي لمداخله؛ وآلية تقييد نمو الأوزان في المشبك لم تُحدَّد. & \cite{Oja1982} & نظرية \\
7 & هيب بالتوقيت: نافذة نحو $\pm20\ms$ (الشكل~\ref{fig:stdp})؛ ومن التسلسلات المتكررة تنمو السلاسل & أزواج من عصبونات الحُصين في المزرعة: نبضة المتلقي خلال $20\ms$ بعد نبضة المرسِل تعطي تقوية ثابتة، وخلال $20\ms$ قبلها تعطي إضعافًا؛ وكلاهما يعتمد على مستقبلات \foreignlanguage{english}{NMDA}. النسبة $A_-=0.6A_+$ في الشكل توضيحية. وأن السلاسل تنمو هكذا في الشبكة لم يُقَس مباشرة. & \cite{BiPoo1998} & مُقاس \\
8 & الأثر~\eqref{eq:threefactor}: الدوبامين الواصل بعد $1\text{--}2\,\text{s}$ من النشاط المشترك يقوّي الوصلة، وبعد ذلك لا (القضية~\ref{prop:threefactor}) & شرائح المخطط، تحفيز بصري جيني منفصل لمداخل الغلوتامات والدوبامين: لا تنمو الشويكة إلا إذا وصل الدوبامين بعد $0.3\text{--}2\,\text{s}$ من دخل الغلوتامات؛ والنافذة يحددها ارتفاع \foreignlanguage{english}{cAMP} السريع وتفككه. وفي القشرة يترك الإقران آثارًا منفصلة للتقوية والإضعاف، تحوّلها مستقبلات الأمينات الأحادية إلى تغيرات طويلة الأمد في نافذة من رتبة الثواني. & \cite{Yagishita2014, He2015eligibility} & مُقاس \\
9 & توجد نوافذ طولها ثوانٍ بلا دوبامين أيضًا: الإشارة الثالثة إثارة قوية للمتلقي & فأر حي، المنطقة \foreignlanguage{english}{CA1}: حدث هضبة واحد في المتلقي يقوّي المداخل الواصلة قبله وبعده بثوانٍ، وينشئ في عبور واحد حقلَ مكان. & \cite{Bittner2017} & مُقاس \\
10 & الخطوة المتوسطة لقاعدة العوامل الثلاثة تسير على تدرج المكافأة~\eqref{eq:perturbation} (القضية~\ref{prop:gradient}) & مبرهنة التدرج للتعلّم بالانحرافات العشوائية؛ واشتُقت في الفصل لضجيج صغير. & \cite{Williams1992} & نظرية \\
11 & الضجيج بحث: تتعلم الشبكة من انحرافاتها العشوائية & طيور الحسون الصغيرة: تعطيل النواة التي تخرج من دارة العقد القاعدية يزيل تنوّع الأغنية بحدة وبصورة عكوسة. الحسون البالغ: يُطبَّق تشويش على أداءات المقطع التي تقع طبقتها على جانب من الوسط، فتزيح الطيور الطبقة بسرعة إلى الجانب الآخر، أي تتعلم من تشتتها هي. & \cite{Olveczky2005, TumerBrainard2007} & مُقاس \\
12 & الاختيار بين اثنين يُتعلَّم عند $\tau_e=1\,\text{s}$ و$\delta=R-\bar R$؛ ولا يُتعلَّم عند $\tau_e=50\ms$؛ وبلا طرح تنمو الأوزان بلا حد، ومع معدل تعلّم كبير قد تثبّت الشبكة الاختيار الأسوأ & \texttt{models/dofa\_learning.py}، $\eta=0.05$، $500$ محاولة، $6$ بذور. $\tau_e=1\,\text{s}$: نسبة اختيار $A$ $0.65\text{--}0.79\to0.96\text{--}1.00$، والأوزان $0.85\text{--}1.33$. $\tau_e=50\ms$: $0.38\text{--}0.55$، والأوزان لا تتغير. $\delta=R$: وزن الفائز $860\text{--}1190$. $\delta=R$، $\eta=0.5$، $3$ بذور: واحدة ثبتت على $B$ ($0.04\to0$). يطبع البرنامج الحالات الأربع كلها؛ وإعادة الحساب طابقت الفصل. & \texttt{models/\allowbreak dofa\_\allowbreak learning.py} & نموذج \\
13 & زمن التعلّم بإشارة مشتركة واحدة يزداد متناسبًا مع عدد العصبونات الضاجّة (القضية~\ref{prop:broadcastcost}) & منحنيات التعلّم لشبكة خطية: اضطراب العصبونات أبطأ من التدرج الدقيق بعدد مرات يساوي عدد عصبونات الخرج، واضطراب الأوزان أبطأ فوق ذلك بعدد مرات يساوي عدد المداخل. & \cite{Werfel2005} & نظرية \\
14 & مخارج \nt{DOPA} تذهب قبل كل شيء إلى مناطق اختيار الأفعال & إعادة بناء كاملة لمحاوير عصبونات \nt{DOPA} سوداوية مخططية مفردة عند الجرذ: فروع الخلية الواحدة تغطي $0.45\text{--}5.7\,\%$ من حجم المخطط (بمتوسط $2.7\,\%$)، ولا تكاد توجد فروع خارج المخطط. & \cite{Matsuda2009} & مُقاس \\
15 & تتعلم القشرة التنبؤ بدخلها، ويُحسب الخطأ في مكانه & مراجعة: في القشرة الحسية استجابات لعدم التطابق بين الدخل المتوقع والواصل. أما أن هذا قاعدة تعلّم عامة للقشرة فغير مبرهن. & \cite{Keller2018} & فرضية \\
16 & يسلك الدوبامين سلوك الخطأ~\eqref{eq:ctd}: ذروة، وانتقال إلى الإشارة المنبئة، وهبوط (الشكل~\ref{fig:ctdcases}) & عصبونات \nt{DOPA} عند القرد: ذروة للعصير غير المتوقع؛ وبعد التعلّم ذروة للإشارة المنبئة ولا استجابة للعصير الموعود؛ وهبوط حين لا يصل العصير. الصيغة المتصلة~\eqref{eq:ctd} نظرية. & \cite{Schultz1997, Doya2000} & مُقاس \\
17 & دارة تحسب $\dd V/\dd t$ وتطرح التوقع & فئران، تسجيل وعلم بصري جيني في السقيفة البطنية: عصبونات \nt{DOPA} تطرح التوقع من المكافأة؛ وعصبونات \nt{GABA} المجاورة تثبّطها حين تكون المكافأة متوقعة، وإثارتها أو تثبيطها يغيّر الخطأ. أما كيف تُحسب المشتقة فلم يُبيَّن. & \cite{Eshel2015arithmetic} & فرضية \\
18 & عصبون \nt{DOPA} ناظمة إيقاع؛ والهبوطات أضحل من الذرى & في الشريحة تفرّغ عصبونات \nt{DOPA} تلقائيًا وبانتظام شديد، على كمون ناظم بطيء. وعند القرد يتبع التردد الخطأ الموجب كميًا، أما السالب فلا ينقله إلا ضعيفًا. & \cite{GraceOnn1989, Bayer2005} & مُقاس \\
19 & الفاعل والناقد (الشكل~\ref{fig:actorcritic}) & الخوارزمية من نظرية التعلّم المعزَّز. عند البشر (\foreignlanguage{english}{fMRI}) يسلك المخطط البطني سلوك الناقد في الاختيار ودونه، أما الظهري فلا يفعل ذلك إلا حين يجب اختيار الأفعال. & \cite{SuttonBarto2018, ODoherty2004actor} & فرضية \\
20 & إشارة $\delta$ في القاعدة معكوسة في العصبونات ذات العائلة الثانية من المستقبلات & شرائح المخطط: في العصبونات ذات مستقبلات \foreignlanguage{english}{D1} يعطي الإقران تقوية تحتاج إلى \foreignlanguage{english}{D1}؛ وفي ذات \foreignlanguage{english}{D2} إضعافًا يحتاج إلى \foreignlanguage{english}{D2}. & \cite{Shen2008} & مُقاس \\
\bottomrule
\end{longtable}}
